\scanpage{221}

\mspage{95}

(7.4.) e$'$) les \uline{$\lim$} filtrantes dans $M$ commutent au chgt de base, il suffit \unsure{=|même} de le supposer dans le cas très particulier \unsure{qu'on a} envisagé ici. \unsure{//}

Comme \deleted{$C$ \inserted{\deleted{$=\mathrm{Mon}$}} et $W$ sont $\Pi$-} $W$ est $\Pi$-accessible, et $(f_i)_{i\in I}$ est $\Pi$-adapté \deleted{\unsure{avec}}, comme $\lim f_i\in W$, il s'ensuit que l'ens \deleted{des} $I_W$ des $i\in I$ tels que $f_i\in W$ est cofinal. Comme $A\overset{f}{\hookrightarrow}B$ n'est pas iso, il existe donc un $i\in I_W$ tel que $B_i\hookrightarrow B$ ne se factorise pas par $A$. On a donc le diagramme

\diagram{\begin{tikzcd}[ampersand replacement=\&,column sep=small,row sep=small]
 A_i \arrow[r,hook,dash,"f_i"] \arrow[d,hook] \arrow[dr,phantom,"\mathrm{cocart}"] \& B_i \arrow[d,hook] \arrow[dr,dashed,"\mathrm{incl.}"] \& {}\\
 A \arrow[r,hook,"\alpha_i"] \arrow[rr,bend right=19,dashed,"f"'] \& C_i \arrow[r,hook,"\beta_i"] \& B
\end{tikzcd}}.

\deleted{On} avec $f_i\in(C\cap W)^\Pi\inserted{=C_0}$, donc $\alpha_i\in\widetilde C_0$; \deleted{\unsure{et}} On \deleted{\illegible{3}} \inserted{a besoin} de savoir que $\alpha_i\in W$, donc une condition supplémentaire sur $W$ :

(7.5.) Si on a un carré cocartésien

\diagram{\begin{tikzcd}[ampersand replacement=\&,column sep=small,row sep=small]
 A \arrow[r,hook,"i"] \arrow[d] \& B \arrow[d]\\
 A' \arrow[r,hook,"i'"] \& B'
\end{tikzcd}}\hfill,

avec $i\in\mathrm{Mon}\cap W$, alors $i'\in W$

% slanted
\margin{\uline{NB} Cette condition est vérifiée}

(j'ai renforcé (\unsure{7|4}.1.4.) en conséquence)\unsure{.}
% archive sheet number: 220
